%%%%%%%%%%%%%%%%%%%%%%%%%%%%%%%%%%%%%%%%%%%%%%%%%%%%%%%%%%%%%%%%%%%%%%%%%%%%%%%%
%% 15-PHASE FRAMEWORK: VISUAL CHARTS AND INFOGRAPHICS
%% TikZ-only visuals (NO PNG, NO \includegraphics)
%% Uses pgfplots for bar charts, raw TikZ for flow diagrams
%%
%% Required packages (already loaded in main.tex):
%%   tikz, pgfplots (compat=1.18), needspace, xcolor, caption
%%
%% Color: ggunavy (#003057) already defined in main.tex
%%
%% Charts provided as \newcommand macros for flexible inclusion:
%%   \phaseChapterCoverage   — Horizontal bar chart: phases per chapter
%%   \taxonomyPassRates      — Grouped bar chart: 525-module pass rates
%%   \diseaseAccuracyChart   — Horizontal bar chart: ASD classification accuracy
%%   \evidenceTraceability   — Sankey-style flow: Problems → Contributions
%%%%%%%%%%%%%%%%%%%%%%%%%%%%%%%%%%%%%%%%%%%%%%%%%%%%%%%%%%%%%%%%%%%%%%%%%%%%%%%%

\normalsize\normalfont\color{black}


%% ═══════════════════════════════════════════════════════════════════════════════
%% CHART 1: PHASE COVERAGE BY CHAPTER (Horizontal Bar Chart)
%% Shows how many of the 15 research phases each thesis chapter covers.
%% Ch.1: 1 phase, Ch.2: 3 phases, Ch.3: 4 phases, Ch.4: 2 phases, Ch.5: 5 phases
%% ═══════════════════════════════════════════════════════════════════════════════
\newcommand{\phaseChapterCoverage}{%
\noindent\textbf{Distribution of Research Phases Across Thesis.}

\needspace{20\baselineskip}%
\begin{figure}[!htbp]
\centering
\resizebox{0.95\textwidth}{!}{%
\begin{tikzpicture}
\begin{axis}[
    xbar,
    width=0.78\textwidth,
    height=8cm,
    bar width=14pt,
    xmin=0, xmax=7,
    xtick={0,1,2,3,4,5,6,7},
    xlabel={\small Number of Research Phases},
    xlabel style={font=\small},
    ytick=data,
    yticklabels={%
        {\small Ch.~5: Discussion \&}\\{\small Recommendations},
        {\small Ch.~4: Data Analysis}\\{\small \& Findings},
        {\small Ch.~3: Research}\\{\small Methods},
        {\small Ch.~2: Literature}\\{\small Review},
        {\small Ch.~1: Introduction}%
    },
    yticklabel style={align=right, text width=3.8cm, anchor=east},
    xmajorgrids=true,
    grid style={gray!25, dashed},
    axis line style={ggunavy, line width=0.8pt},
    tick style={ggunavy},
    tick label style={color=black},
    enlarge y limits=0.15,
    nodes near coords,
    nodes near coords style={font=\small\bfseries, color=white, anchor=east, xshift=-8pt},
    every axis plot/.append style={fill=ggunavy, draw=ggunavy!80!black},
    clip=false,
]
\addplot coordinates {
    (1,0)
    (3,1)
    (4,2)
    (2,3)
    (5,4)
};
%% Annotation: total phases
\node[anchor=west, font=\small\color{ggunavy}] at (axis cs:5.3,4)
    {2 core + 3 meta};
\end{axis}
\end{tikzpicture}%
}% end resizebox
\caption[Distribution of Research Phases Across Thesis Chapters]{Distribution of Research Phases Across Thesis Chapters. Each of the fifteen doctoral research phases is mapped to a primary thesis chapter. Chapter~5 covers the largest share, including two core interpretation phases (Phases~11--12) and three cross-cutting meta-phases (Phases~13--15). Chapter~3 covers four consecutive design phases (Phases~5--8).}


\label{fig:phase_chapter_distribution}
\end{figure}
}


%% ═══════════════════════════════════════════════════════════════════════════════
%% CHART 2: QUALITY TAXONOMY PASS RATES (Grouped Bar Chart)
%% 525-Module Unified Quality Taxonomy across three tiers.
%% Tier 1: 8-Phase DL (111 modules, 97.4%)
%% Tier 2: 13-Phase GenAI QA (220 modules, 98.1%)
%% Tier 3: 10-Pillar Governance (194 modules, 98.4%)
%% ═══════════════════════════════════════════════════════════════════════════════
\newcommand{\taxonomyPassRates}{%
\noindent\textbf{525-Module Unified Quality Taxonomy: Pass Rates.}

\needspace{20\baselineskip}%
\begin{figure}[!htbp]
\centering
\resizebox{0.95\textwidth}{!}{%
\begin{tikzpicture}
\begin{axis}[
    ybar,
    width=0.82\textwidth,
    height=9cm,
    bar width=30pt,
    ymin=90, ymax=102,
    ytick={90,91,92,93,94,95,96,97,98,99,100},
    ylabel={\small Pass Rate (\%)},
    ylabel style={font=\small},
    ymajorgrids=true,
    grid style={gray!25, dashed},
    axis line style={ggunavy, line width=0.8pt},
    tick style={ggunavy},
    tick label style={color=black, font=\small},
    xtick=data,
    xticklabels={%
        {\small Tier~1:}\\{\small 8-Phase DL},
        {\small Tier~2:}\\{\small 13-Phase GenAI QA},
        {\small Tier~3:}\\{\small 10-Pillar Governance}%
    },
    xticklabel style={align=center, text width=3.5cm, anchor=north},
    enlarge x limits=0.25,
    nodes near coords={\pgfmathprintnumber\pgfplotspointmeta\%},
    nodes near coords style={font=\small\bfseries, color=ggunavy, anchor=south, yshift=2pt},
    clip=false,
    legend style={at={(0.5,-0.22)}, anchor=north, font=\small,
                  draw=ggunavy, fill=white, legend columns=1},
]
%% Pass rate bars with different opacities
\addplot[fill=ggunavy, draw=ggunavy!80!black] coordinates {(1,97.4)};
\addplot[fill=ggunavy!75, draw=ggunavy!60!black] coordinates {(2,98.1)};
\addplot[fill=ggunavy!50, draw=ggunavy!40!black] coordinates {(3,98.4)};

%% Module count annotations below each bar
\node[font=\small\color{ggunavy!70}, anchor=north] at (axis cs:1,90.6)
    {111 modules};
\node[font=\small\color{ggunavy!70}, anchor=north] at (axis cs:2,90.6)
    {220 modules};
\node[font=\small\color{ggunavy!70}, anchor=north] at (axis cs:3,90.6)
    {194 modules};

%% Total annotation
\node[anchor=west, font=\small\bfseries\color{ggunavy},
      draw=ggunavy!30, fill=ggunavy!5, rounded corners=2pt, inner sep=4pt]
    at (axis cs:2.8,99.2) {525 modules total};

\end{axis}
\end{tikzpicture}%
}% end resizebox
\caption[525-Module Unified Quality Taxonomy]{525-Module Unified Quality Taxonomy: Pass Rates by Tier. The three-tier quality taxonomy validates all research artefacts against 525 discrete quality modules. Tier~1 (8-Phase DL Analysis) evaluates 111 modules at 97.4\% pass rate; Tier~2 (13-Phase GenAI QA) evaluates 220 modules at 98.1\%; Tier~3 (10 Cross-Cutting AI Governance Pillars) evaluates 194 modules at 98.4\%. The progressive improvement across tiers reflects iterative quality refinement.}


\label{fig:taxonomy_pass_rates}
\end{figure}
}


%% ═══════════════════════════════════════════════════════════════════════════════
%% CHART 3: DISEASE DOMAIN ACCURACY (Horizontal Bar Chart)
%% Per-domain classification accuracy sorted descending, with ASD ensemble baseline.
%% PD: 100%, Epilepsy: 99.02%, ASD: 97.67%, Schizophrenia: 97.17%,
%% Stress: 94.17%, Depression: 91.07%, ASD ensemble consensus: 97.67%
%% ═══════════════════════════════════════════════════════════════════════════════
\newcommand{\diseaseAccuracyChart}{%
\noindent\textbf{Per-Domain Classification Accuracy and ASD.}

\needspace{20\baselineskip}%
\begin{figure}[!htbp]
\centering
\resizebox{0.95\textwidth}{!}{%
\begin{tikzpicture}
\begin{axis}[
    xbar,
    width=0.84\textwidth,
    height=10cm,
    bar width=13pt,
    xmin=85, xmax=102,
    xtick={86,88,90,92,94,96,98,100},
    xlabel={\small Classification Accuracy (\%)},
    xlabel style={font=\small},
    ytick=data,
    yticklabels={%
        {\small Depression},
        {\small Stress},
        {\small Schizophrenia},
        {\small ASD (Autism)},
        {\small Epilepsy},
        {\small PD (Parkinson)}%
    },
    yticklabel style={align=right, text width=3cm, anchor=east, font=\small},
    xmajorgrids=true,
    grid style={gray!25, dashed},
    axis line style={ggunavy, line width=0.8pt},
    tick style={ggunavy},
    tick label style={color=black},
    enlarge y limits=0.12,
    nodes near coords={\pgfmathprintnumber[fixed, precision=2]\pgfplotspointmeta\%},
    nodes near coords style={font=\small\bfseries, color=ggunavy, anchor=west, xshift=3pt},
    every axis plot/.append style={fill=ggunavy, draw=ggunavy!80!black},
    clip=false,
]
%% Bars sorted by accuracy ascending (bottom to top)
\addplot coordinates {
    (91.07,0)
    (94.17,1)
    (97.17,2)
    (97.67,3)
    (99.02,4)
    (100.00,5)
};

%% ASD ensemble consensus baseline — dashed vertical reference line
\draw[dashed, line width=1.2pt, ggunavy!60]
    (axis cs:96.4,\pgfkeysvalueof{/pgfplots/ymin}-0.6) --
    (axis cs:96.4,\pgfkeysvalueof{/pgfplots/ymax}+0.5);
\node[anchor=south west, font=\small\color{ggunavy!70}]
    at (axis cs:96.4,5.3) {ASD Accuracy: 97.67\%};

\end{axis}
\end{tikzpicture}%
}% end resizebox
\caption[Per-Domain Classification Accuracy and ASD ensemble Consensus Baseline]{Per-Domain Classification Accuracy and ASD ensemble Consensus Baseline. Individual disease domains achieve 91--100\% classification accuracy using optimised DL pipelines. The dashed vertical line marks the Multi-System Consensus Architecture (ASD ensemble) ASD-focused consensus accuracy of 97.67\%~$\pm$~2.1\%, which represents the consensus when integrating predictions for ASD simultaneously. PD achieves perfect classification; Depression is the most challenging domain.}


\label{fig:disease_accuracy}
\end{figure}
}


%% ═══════════════════════════════════════════════════════════════════════════════
%% CHART 4: RESEARCH EVIDENCE TRACEABILITY CHAIN (Sankey-style Flow)
%% Left-to-right flow: Problems (P1-P4) → Gaps (G1-G12) → Questions (RQ1-RQ6)
%% → Hypotheses (H1-H5) → Contributions (C1-C8)
%% ═══════════════════════════════════════════════════════════════════════════════
\newcommand{\evidenceTraceability}{%
\noindent\textbf{Research Evidence Traceability Chain: From.}

\needspace{24\baselineskip}%
\begin{figure}[!htbp]
\centering
\resizebox{0.95\textwidth}{!}{%
\begin{tikzpicture}[
    %% Node styles — color-coded by stage
    probnode/.style={rectangle, draw=ggunavy, fill=ggunavy, text=white,
                     text width=2.4cm, minimum height=0.72cm, align=center,
                     font=\small\bfseries, rounded corners=2pt,
                     line width=0.6pt},
    gapnode/.style={rectangle, draw=ggunavy!75!black, fill=ggunavy!75,
                    text=white, text width=1cm, minimum height=0.52cm,
                    align=center, font=\small\bfseries,
                    rounded corners=2pt, line width=0.5pt},
    rqnode/.style={rectangle, draw=ggunavy!55!black, fill=ggunavy!55,
                   text=white, text width=1.3cm, minimum height=0.62cm,
                   align=center, font=\small\bfseries,
                   rounded corners=2pt, line width=0.5pt},
    hypnode/.style={rectangle, draw=ggunavy!40!black, fill=ggunavy!40,
                    text=white, text width=1.2cm, minimum height=0.62cm,
                    align=center, font=\small\bfseries,
                    rounded corners=2pt, line width=0.5pt},
    contnode/.style={rectangle, draw=ggunavy!25!black, fill=ggunavy!25,
                     text width=2.4cm, minimum height=0.72cm, align=center,
                     font=\small\bfseries, rounded corners=2pt,
                     line width=0.5pt},
    stageheader/.style={font=\small\bfseries\color{ggunavy}, anchor=south},
    flowarrow/.style={-{Stealth[length=3.5mm, width=2.5mm]},
                      line width=0.8pt, ggunavy!50},
    %% Bundle arrow for many-to-many connections
    bundlearrow/.style={-{Stealth[length=3.5mm, width=2.5mm]},
                        line width=0.5pt, ggunavy!35},
]

%% ── Column positions ──
\def\colP{0}        % Problems
\def\colG{4.2}      % Gaps
\def\colR{7.5}      % Research Questions
\def\colH{10.8}     % Hypotheses
\def\colC{14.4}     % Contributions

%% ── Stage headers ──
\node[stageheader] at (\colP, 5.0) {Problems};
\node[stageheader] at (\colG, 5.0) {Gaps};
\node[stageheader] at (\colR, 5.0) {Questions};
\node[stageheader] at (\colH, 5.0) {Hypotheses};
\node[stageheader] at (\colC, 5.0) {Contributions};

%% ── Count labels ──
\node[font=\small\color{ggunavy!60}] at (\colP, 4.5) {4 identified};
\node[font=\small\color{ggunavy!60}] at (\colG, 4.5) {12 mapped};
\node[font=\small\color{ggunavy!60}] at (\colR, 4.5) {6 formulated};
\node[font=\small\color{ggunavy!60}] at (\colH, 4.5) {5 tested};
\node[font=\small\color{ggunavy!60}] at (\colC, 4.5) {8 delivered};

%% ── Problems (P1--P4) ──
\node[probnode] (P1) at (\colP, 3.6) {P1: Diagnostic};
\node[probnode] (P2) at (\colP, 2.5) {P2: Integration};
\node[probnode] (P3) at (\colP, 1.4) {P3: Governance};
\node[probnode] (P4) at (\colP, 0.3) {P4: Scalability};

%% ── Gaps (G1--G12, stacked vertically in groups) ──
\node[gapnode] (G1)  at (\colG, 4.0)  {G1};
\node[gapnode] (G2)  at (\colG, 3.4)  {G2};
\node[gapnode] (G3)  at (\colG, 2.8)  {G3};
\node[gapnode] (G4)  at ($(\colG,0)+(0.9, 4.0)$) {G4};
\node[gapnode] (G5)  at ($(\colG,0)+(0.9, 3.4)$) {G5};
\node[gapnode] (G6)  at ($(\colG,0)+(0.9, 2.8)$) {G6};
\node[gapnode] (G7)  at (\colG, 2.0)  {G7};
\node[gapnode] (G8)  at ($(\colG,0)+(0.9, 2.0)$) {G8};
\node[gapnode] (G9)  at (\colG, 1.2)  {G9};
\node[gapnode] (G10) at ($(\colG,0)+(0.9, 1.2)$) {G10};
\node[gapnode] (G11) at (\colG, 0.4)  {G11};
\node[gapnode] (G12) at ($(\colG,0)+(0.9, 0.4)$) {G12};

%% ── Research Questions (RQ1--RQ8) ──
\node[rqnode] (RQ1) at (\colR, 3.8) {RQ1};
\node[rqnode] (RQ2) at (\colR, 3.0) {RQ2};
\node[rqnode] (RQ3) at (\colR, 2.2) {RQ3};
\node[rqnode] (RQ4) at (\colR, 1.4) {RQ4};
\node[rqnode] (RQ5) at (\colR, 0.6) {RQ5};
\node[rqnode] (RQ6) at (\colR, -0.2) {RQ6};

%% ── Hypotheses (H1--H5) ──
\node[hypnode] (H1) at (\colH, 3.6) {H1};
\node[hypnode] (H2) at (\colH, 2.7) {H2};
\node[hypnode] (H3) at (\colH, 1.8) {H3};
\node[hypnode] (H4) at (\colH, 0.9) {H4};
\node[hypnode] (H5) at (\colH, 0.0) {H5};

%% ── Contributions (C1--C8) ──
\node[contnode] (C1) at (\colC, 4.0)  {C1: RGAIG};
\node[contnode] (C2) at (\colC, 3.2)  {C2: ASD Ensemble};
\node[contnode] (C3) at (\colC, 2.4)  {C3: DL};
\node[contnode] (C4) at (\colC, 1.6)  {C4: GenAI QA};
\node[contnode] (C5) at (\colC, 0.8)  {C5: Governance};
\node[contnode] (C6) at (\colC, 0.0)  {C6: Wearable};
\node[contnode] (C7) at (\colC, -0.8) {C7: RAG};
\node[contnode] (C8) at (\colC, -1.6) {C8: ROI Model};

%% ══════════════════════════════════════════════════════════════════════
%% ARROWS: Problems → Gaps (bundle arrows to reduce visual clutter)
%% ══════════════════════════════════════════════════════════════════════
%% P1 → G1, G2, G3
\draw[bundlearrow] (P1.east) -- (G1.west);
\draw[bundlearrow] (P1.east) -- (G2.west);
\draw[bundlearrow] (P1.east) -- (G3.west);
%% P2 → G4, G5, G6, G7
\draw[bundlearrow] (P2.east) -- (G4.west);
\draw[bundlearrow] (P2.east) -- (G5.west);
\draw[bundlearrow] (P2.east) -- (G6.west);
\draw[bundlearrow] (P2.east) -- (G7.west);
%% P3 → G8, G9, G10
\draw[bundlearrow] (P3.east) -- (G8.west);
\draw[bundlearrow] (P3.east) -- (G9.west);
\draw[bundlearrow] (P3.east) -- (G10.west);
%% P4 → G11, G12
\draw[bundlearrow] (P4.east) -- (G11.west);
\draw[bundlearrow] (P4.east) -- (G12.west);

%% ══════════════════════════════════════════════════════════════════════
%% ARROWS: Gaps → Research Questions (bundled)
%% ══════════════════════════════════════════════════════════════════════
%% G1--G3 → RQ1, RQ2
\draw[bundlearrow] (G3.east) -- (RQ1.west);
\draw[bundlearrow] (G3.east) -- (RQ2.west);
%% G4--G6 → RQ2, RQ3
\draw[bundlearrow] (G6.east) -- (RQ2.west);
\draw[bundlearrow] (G6.east) -- (RQ3.west);
%% G7, G8 → RQ3, RQ4
\draw[bundlearrow] (G8.east) -- (RQ3.west);
\draw[bundlearrow] (G8.east) -- (RQ4.west);
%% G9, G10 → RQ4, RQ5
\draw[bundlearrow] (G10.east) -- (RQ4.west);
\draw[bundlearrow] (G10.east) -- (RQ5.west);
%% G11, G12 → RQ5, RQ6
\draw[bundlearrow] (G12.east) -- (RQ5.west);
\draw[bundlearrow] (G12.east) -- (RQ6.west);

%% ══════════════════════════════════════════════════════════════════════
%% ARROWS: Research Questions → Hypotheses
%% ══════════════════════════════════════════════════════════════════════
\draw[flowarrow] (RQ1.east) -- (H1.west);
\draw[flowarrow] (RQ2.east) -- (H1.west);
\draw[flowarrow] (RQ2.east) -- (H2.west);
\draw[flowarrow] (RQ3.east) -- (H2.west);
\draw[flowarrow] (RQ3.east) -- (H3.west);
\draw[flowarrow] (RQ4.east) -- (H3.west);
\draw[flowarrow] (RQ4.east) -- (H4.west);
\draw[flowarrow] (RQ5.east) -- (H4.west);
\draw[flowarrow] (RQ5.east) -- (H5.west);
\draw[flowarrow] (RQ6.east) -- (H5.west);

%% ══════════════════════════════════════════════════════════════════════
%% ARROWS: Hypotheses → Contributions
%% ══════════════════════════════════════════════════════════════════════
\draw[flowarrow] (H1.east) -- (C1.west);
\draw[flowarrow] (H1.east) -- (C2.west);
\draw[flowarrow] (H2.east) -- (C2.west);
\draw[flowarrow] (H2.east) -- (C3.west);
\draw[flowarrow] (H3.east) -- (C3.west);
\draw[flowarrow] (H3.east) -- (C4.west);
\draw[flowarrow] (H4.east) -- (C5.west);
\draw[flowarrow] (H4.east) -- (C6.west);
\draw[flowarrow] (H5.east) -- (C6.west);
\draw[flowarrow] (H5.east) -- (C7.west);
\draw[flowarrow] (H5.east) -- (C8.west);

%% ══════════════════════════════════════════════════════════════════════
%% Stage divider lines (subtle vertical separators)
%% ══════════════════════════════════════════════════════════════════════
\foreach \x in {2.6, 6.4, 9.2, 12.6} {
    \draw[ggunavy!15, line width=0.4pt, dashed]
        (\x, -2.0) -- (\x, 4.8);
}

%% ── Bottom legend: color gradient explanation ──
\node[anchor=north west, font=\small\color{ggunavy}] at (0, -2.3)
    {Colour intensity decreases left to right, reflecting the
     abstraction flow from concrete problems to generalised contributions.};

\end{tikzpicture}%
}% end resizebox
\caption[Research Evidence Traceability Chain]{Research Evidence Traceability Chain: From Problems to Contributions. Four research problems (P1--P4) generate twelve literature gaps (G1--G12), which are distilled into eight research questions (RQ1--RQ8), operationalised as five testable hypotheses (H1--H5), and resolved through eight novel contributions (C1--C8). Arrows indicate the logical derivation path; every contribution traces back to at least one empirically tested hypothesis.}


\label{fig:evidence_traceability}
\end{figure}
}


%% ═══════════════════════════════════════════════════════════════════════════════
%% END OF PHASE CHARTS
%% ═══════════════════════════════════════════════════════════════════════════════
